%%% FIELD proof-observation-group-partition
By \(\text{lem:yang-algorithm-equivalence}\), under \(\text{ass:canonical-lvsemme}\), \(\text{ass:minimality}\), \(\text{ass:rank-faithfulness}\), and \(\text{ass:known-observation-labels}\), Algorithm 1 returns an ancestral ordered grouping that is a partition in causal group order. Every substantive label is assigned to exactly one recovered observation group. Restricting admissible center roles by the known map \(\ell\) can delete a center candidacy, but cannot duplicate a label or assign it to a second group. Hence the nonempty candidate-center sets \(R_1,\ldots,R_k\) are pairwise disjoint.

By the definition of \(V_{\mathrm{fix}}\),
\[
V_{\mathrm{fix}}(\mathcal G,\ell)=V\setminus\bigcup_{g=1}^{k}R_g.
\]
Pairwise disjointness therefore gives
\[
V=\left(\bigsqcup_{g=1}^{k}R_g\right)\sqcup V_{\mathrm{fix}}(\mathcal G,\ell).
\]
Thus, for each \(u\in V\), exactly one clause in the definition of \(g_{\mathcal G,\ell}\) applies: either \(u\in R_g\) for a unique \(g\in[k]\), or \(u\in V_{\mathrm{fix}}\) and \(g_{\mathcal G,\ell}(u)=\bot\).

A legal completion \(\omega=((r_h,c_h))_{h=1}^{k}\) selects exactly one row \(r_g\in R_g\). Consequently, if \(u\in R_g\), then \(u\) is the cogent center exactly when \(r_g=u\); every other row of that recovered observation group has the mleaf role. By the label and singleton constraints in the definition of \(\Omega\), every label in \(V_{\mathrm{fix}}\) remains a fixed singleton mleaf. These conclusions use the corrected Yang minimality and rank-faithfulness assumptions only through the cited algorithm-equivalence result.

%%% FIELD proof-normalized-reconstruction
Fix \(\omega\in\Omega(W,\mathcal G,\ell)\), and abbreviate \(R=R(\omega)\), \(C=C(\omega)\), \(L=L(\omega)\), and \(H=H(\omega)\). The pivot and causal-order clauses in legality give
\[
W_{r_gc_g}\ne0\quad(g\in[k]),\qquad W_{r_gc_h}=0\quad(g<h).
\]
It follows that \(\Delta_\omega\) is invertible and
\[
(Q_\omega)_{gg}=1,\qquad (Q_\omega)_{gh}=0\quad(g<h).
\]
Thus \(Q_\omega\) is unit lower triangular and invertible.

Let \(N_C\) denote the selected independent structural sources and put \(\overline N_C=\Delta_\omega N_C\). In substantive row order \((R,L)\), the reconstructed equations are
\[
X_R=A_\omega X_R+B^C_\omega N_H+\overline N_C,
\qquad X_L=D_\omega X_R+B^L_\omega N_H.
\]
Because \(I-A_\omega=Q_\omega^{-1}\), their retained mixing matrix in source order \((\overline N_C,N_H)\) is
\[
\overline W_\omega=
\begin{pmatrix}
Q_\omega&Q_\omega B^C_\omega\\
D_\omega Q_\omega&D_\omega Q_\omega B^C_\omega+B^L_\omega
\end{pmatrix}.
\]
The defining identities give
\[
Q_\omega=W[R,C]\Delta_\omega^{-1},\qquad Q_\omega B^C_\omega=W[R,H],
\]
\[
D_\omega Q_\omega=W[L,C]\Delta_\omega^{-1},\qquad
D_\omega Q_\omega B^C_\omega+B^L_\omega=W[L,H].
\]
Therefore
\[
\overline W_\omega=
\begin{pmatrix}
W[R,C]\Delta_\omega^{-1}&W[R,H]\\
W[L,C]\Delta_\omega^{-1}&W[L,H]
\end{pmatrix}.
\]
Returning from \(\overline N_C\) to \(N_C\) multiplies the selected columns by \(\Delta_\omega\), so the retained mixing matrix is exactly \(W\).

The inverse of a unit lower-triangular matrix is unit lower triangular. Hence \(A_\omega=I-Q_\omega^{-1}\) is strictly lower triangular, proving acyclicity of the center submodel. The block \(D_\omega\) points from centers to mleaf rows, \(B^C_\omega,B^L_\omega\) load root latent sources, and the right substantive blocks are zero because mleafs have no substantive children. All remaining observation-label, latent-child, pivot, and canonical support conditions are clauses of \(\omega\in\Omega\). Thus \(M_\omega\) is legal and canonical.

Let \(S\) be an invertible diagonal column-scaling matrix and put \(W'=WS\). For the selected and complementary diagonal blocks,
\[
\Delta'_\omega=\Delta_\omega S_C,
\qquad W'[R,C](\Delta'_\omega)^{-1}=W[R,C]\Delta_\omega^{-1}=Q_\omega.
\]
Thus \(A'_\omega=A_\omega\) and \(D'_\omega=D_\omega\), whereas
\[
(B^C_\omega)'=B^C_\omega S_H,
\qquad (B^L_\omega)'=B^L_\omega S_H.
\]
The latter changes only the units of independent latent sources. Column permutation merely relabels selected and complementary sources. This is precisely invariance under the recovered permutation-and-scaling indeterminacy.

%%% FIELD proof-local-effect-column
Fix a legal \(\omega\). In substantive row order \((R(\omega),L(\omega))\), the coefficient matrix among substantive variables is
\[
G_\omega=\begin{pmatrix}A_\omega&0\\D_\omega&0\end{pmatrix}.
\]
The zero right-hand blocks are the canonical mleaf condition. Since \(I-A_\omega=Q_\omega^{-1}\), block multiplication gives
\[
(I-G_\omega)^{-1}=\begin{pmatrix}Q_\omega&0\\D_\omega Q_\omega&I\end{pmatrix}.
\]

Acyclicity makes \(G_\omega\) nilpotent, so
\[
(I-G_\omega)^{-1}=I+G_\omega+G_\omega^2+\cdots.
\]
Its \((v,u)\) entry is the sum of products of structural coefficients along all directed paths from \(u\) to \(v\), with the length-zero path contributing one when \(v=u\). Under the surgical intervention fixing \(X_u=x\), differentiating the recursively ordered structural equations with respect to \(x\) gives this same path recursion. Thus the labeled linear total-effect column is the \(u\)-th column of \((I-G_\omega)^{-1}\); this proves its relation to the structural intervention rather than defining it from \(W\).

If \(u\in L(\omega)\), the displayed inverse has \(u\)-th column \(e_u\), so the diagonal effect is one and every off-diagonal effect is zero. By \(\text{lem:observation-group-partition}\), this covers both \(g_{\mathcal G,\ell}(u)=\bot\) and the case \(g_{\mathcal G,\ell}(u)=g\in[k]\) with \(r_g\ne u\).

If \(u=r_g\), the same column is column \(g\) of
\[
\begin{pmatrix}Q_\omega\\D_\omega Q_\omega\end{pmatrix}
=W[V,C(\omega)]\Delta_\omega^{-1},
\]
hence equals
\[
\frac{W_{\bullet c_g}}{W_{uc_g}}.
\]
The denominator is nonzero by the legal-pivot condition. Taking coordinate \(v\) proves every scalar branch. The vector depends only on \((r_g,c_g)\), so choices outside the unique group containing \(u\) cannot alter the entire intervention column.

%%% FIELD proof-complete-effect-image
Fix a tuple covered by the theorem. By \(\text{thm:observational-dog-column-completeness}\), the observationally compatible minimum-edge DOG class is indexed by \(\Omega^\star\). By \(\text{thm:deterministic-joint-effect-image}\), its exact joint \(u\)-intervention-column image is \(S_{\bullet u}\). The definition \(\text{def:effect-image}\) then takes the coordinate projection
\[
S_{vu}=\{x_v:x\in S_{\bullet u}\}.
\]
Consequently
\[
S_{vu}=\{\tau_\omega(v\leftarrow u):\omega\in\Omega^\star\}.
\]
This identifies the causal image with a functional of the recovered mixing class; it does not define the causal effect.

If \(g_{\mathcal G,\ell}(u)=\bot\), the joint set is \(\{e_u\}\), so \(S_{uu}=\{1\}\) and \(S_{vu}=\{0\}\) for \(v\ne u\). On the cogent-group branch, taking coordinate \(v\) gives
\[
S_{vu}=\left\{\frac{W_{vc}}{W_{uc}}:c\in L_u\right\}
\cup\bigl(\{0\}\text{ if }z_u=1\bigr)\quad(v\ne u),
\]
and gives \(S_{uu}=\{1\}\). Every denominator is a nonzero legal pivot. Projection cannot increase cardinality, so
\[
|S_{vu}|\le |S_{\bullet u}|\le q+1.
\]

Every scalar value is a coordinate of a vector in the finite joint image. The joint-image theorem supplies one minimum completion attaining that entire vector. If the scalar image is not a singleton, two such completions attain unequal scalars; observational completeness makes them faithful observationally equivalent countermodels on \(\Theta_{p,q,\ell}^{\mathrm{gen}}\). That theorem also supplies the stratumwise relative-Zariski and Lebesgue-almost-everywhere qualifications. Source-law separability remains separate, and no probability assertion is made for rational input points.

%%% FIELD proof-singleton-minor-corollary
By \(\text{thm:complete-effect-image}\), \(S_{vu}\) is the \(v\)-coordinate projection of the exact joint image. If \(v=u\), every normalized column and \(e_u\) has coordinate one, so \(S_{uu}=\{1\}\). If \(v\ne u\) and \(g_{\mathcal G,\ell}(u)=\bot\), then \(S_{vu}=\{0\}\).

Suppose \(v\ne u\) and \(g_{\mathcal G,\ell}(u)=g\in[k]\). For every \(c\in L_u\), the pair \((u,c)\) occurs in a legal completion, so \(W_{uc}\ne0\). If \(L_u=\varnothing\), nonemptiness of \(\Omega^\star\) gives an attaining pair in group \(g\), necessarily with center different from \(u\). Thus \(z_u=1\), and the image is \(\{0\}\).

If \(z_u=1\), zero belongs to the scalar projection. It is the only value exactly when every displayed ratio is zero. For each \(c\in L_u\),
\[
\frac{W_{vc}}{W_{uc}}=0\quad\Longleftrightarrow\quad W_{vc}=0,
\]
because the pivot is nonzero. This proves the zero branch, including empty \(L_u\) vacuously.

If \(z_u=0\) and \(L_u\ne\varnothing\), no zero is adjoined, and singletonness is equality of all ratios. For \(c,d\in L_u\), both denominators are nonzero, so
\[
\frac{W_{vc}}{W_{uc}}=\frac{W_{vd}}{W_{ud}}
\quad\Longleftrightarrow\quad
W_{vc}W_{ud}-W_{vd}W_{uc}=0.
\]
Thus the pairwise minors are necessary and sufficient, and any admissible ratio is the common value. The cases are exhaustive.

%%% FIELD proof-projected-oracle-reduction
Assume \(\mathsf A_{\min}\) has the stated polynomial bit-time guarantee. One unconstrained call returns
\[
e^\star=\min_{\omega\in\Omega}E(\omega)
\]
and an attaining completion. For each \(g\in[k]\) and \((r,c)\in R_g\times K_g\), make one constrained call and denote its optimum by \(e_{g,r,c}\), with value \(+\infty\) when infeasibility is certified. The constrained feasible set is contained in \(\Omega\), so \(e_{g,r,c}\ge e^\star\). The definition of \(P_g\) gives
\[
(r,c)\in P_g\quad\Longleftrightarrow\quad e_{g,r,c}=e^\star.
\]
For the forward implication, a global minimizer witnessing membership in \(P_g\) is feasible for the constrained problem. For the reverse implication, the returned constrained completion is a global minimizer and witnesses membership. Retaining it gives an attaining witness for every retained pair.

The row sets \(R_g\) are disjoint by \(\text{lem:observation-group-partition}\), and every \(K_g\subseteq C_{\mathrm{src}}\). Hence
\[
\sum_{g=1}^{k}|R_g||K_g|\le q\sum_{g=1}^{k}|R_g|\le pq.
\]
Thus at most \(pq\) constrained calls compute all projections.

For each \(u\), collect retained pairs with first coordinate \(u\), and check whether another first coordinate occurs in its uniquely bound group. The definition \(\text{def:effect-image}\) then forms \(S_{\bullet u}\), and coordinate projection gives every \(S_{vu}\). At most \(q\) normalized columns are processed for each of \(p\) interventions and \(p\) coordinates, requiring \(O(p^2q)\) rational operations. Singletonness uses numerator-zero tests and exact cross-products
\[
W_{vc}W_{ud}=W_{vd}W_{uc},
\]
so certification needs no numerical division.

If an image is a singleton, these comparisons return its common value. If it is not, choose two unequal generated coordinates. A normalized-column value uses the stored minimizer for its pair, whereas an adjoined zero uses a stored alternative-center minimizer. The deterministic theorem proves exact attainment; on the generic population domain, observational completeness also interprets the reconstructions as faithful observationally equivalent countermodels. The bounded oracle calls and rational postprocessing are polynomial under the assumed guarantee.

%%% FIELD proof-no-measurement-reduction
On the stated subclass every substantive variable is directly observed and no substantive row can have an mleaf role. The known-label restriction and \(\text{lem:observation-group-partition}\) give singleton center-row sets \(R_g=\{u_g\}\), no fixed-mleaf labels, and \(z_u=0\) for every \(u\). Thus
\[
S_{\bullet u}(W,\mathcal G(W,\ell),\ell)
=\left\{\frac{W_{\bullet c}}{W_{uc}}:c\in L_u(W,\mathcal G(W,\ell),\ell)\right\},
\]
and
\[
S_{vu}(W,\mathcal G(W,\ell),\ell)
=\left\{\frac{W_{vc}}{W_{uc}}:c\in L_u(W,\mathcal G(W,\ell),\ell)\right\}.
\]
Every denominator is a nonzero legal pivot. With neither measured-variable center exchange nor mleaf rows, the remaining ambiguity is exactly the assignment of an observed structural-noise column versus a root-latent column in canonical latent-variable LiNGAM. Normalizing a candidate column at row \(u\) therefore gives the published normalized mixing-column effect set.

On the common generic domain, \(\text{thm:complete-effect-image}\) makes this coordinate set the complete observational image of the labeled total effect. It is a singleton exactly when that effect is generically identifiable. The revalidated published result \(\text{lem:tramontano-unknown-graph-total-effect-criterion}\) says that, in the same canonical latent-variable LiNGAM class, generic identification holds exactly when its Theorem 3.3 graphical obstruction is absent. Chaining these equivalences proves the proposition. The scalar minor comparison is the algebraic form of the same normalized-column equality test.

%%% FIELD proof-deterministic-joint-effect-image
Fix \(u\in V\). If \(g_{\mathcal G,\ell}(u)=\bot\), the partition lemma makes \(u\) a fixed mleaf, and the local-effect theorem gives
\[
T_{\omega,\bullet u}=e_u\qquad(\omega\in\Omega^\star).
\]
This is the first branch of \(\text{def:effect-image}\).

Suppose \(g_{\mathcal G,\ell}(u)=g\in[k]\). Partition the nonempty finite argmin as
\[
\Omega_u^\star=\{\omega\in\Omega^\star:r_g=u\},
\qquad
\Omega_{\neg u}^\star=\{\omega\in\Omega^\star:r_g\ne u\}.
\]
For \(\omega\in\Omega_u^\star\), the definition of \(P_g\) gives \((u,c_g)\in P_g\), and the local theorem gives
\[
T_{\omega,\bullet u}=\frac{W_{\bullet c_g}}{W_{uc_g}}.
\]
Conversely, if \((u,c)\in P_g\), its defining existential quantifier supplies one \(\omega_c\in\Omega^\star\) with \(\omega_{c,g}=(u,c)\). That one completion attains the whole normalized vector. Its denominator is nonzero because it is a legal pivot.

For \(\omega\in\Omega_{\neg u}^\star\), the partition lemma makes \(u\) an mleaf and the local theorem gives \(T_{\omega,\bullet u}=e_u\). Such a completion exists exactly when some \((r,c)\in P_g\) has \(r\ne u\), namely exactly when \(z_u=1\). These two parts prove the asserted set equality and all attainment claims.

There are at most \(|K_g|\) normalized columns and at most one additional \(e_u\). Since \(K_g\subseteq C_{\mathrm{src}}\),
\[
|S_{\bullet u}|\le |K_g|+1\le q+1.
\]
Taking coordinate \(v\) is exactly the definition of \(S_{vu}\). Finally, each \(P_g\) records only whether a local pair occurs in some global minimizer. Pairs from different groups may have different existential witnesses and need not coexist. Hence no Cartesian factorization of the full total-effect-matrix image follows; it remains indexed by \(\Omega^\star\).

%%% FIELD proof-observational-dog-column-completeness
By \(\text{ass:separability}\), \(P_{\mathrm{obs}}\) identifies the retained mixing matrix up to permutation and nonzero scaling of its columns. By \(\text{lem:normalized-reconstruction}\), the normalized center block, substantive coefficients, and normalized effect columns are invariant under that indeterminacy. By the revalidated \(\text{lem:yang-algorithm-equivalence}\), the corrected assumptions imply that Algorithm 2 enumerates the legal ancestral-ordered-group reconstructions and its minimum-edge output is the Yang DOG class on the Theorem 3 domain. The original model gives one legal choice; finiteness yields \(\Omega^\star\ne\varnothing\).

The simultaneous faithfulness assertion is not inferred from deterministic algebra. The revalidated \(\text{lem:yang-dog-faithfulness-preservation}\) supplies it: the center/noise rewrite to each member of the finite DOG class is invertible, every exceptional Assumption 3 locus has probability zero under the continuous-coefficient regime, and their finite union remains null. The definition of \(\Theta_{p,q,\ell}^{\mathrm{gen}}\) removes those proper algebraic exceptional loci. Hence every attaining reconstruction lies simultaneously in the published rank-faithful class on the stated domain.

Apply \(\text{thm:deterministic-joint-effect-image}\). It gives the exact joint column image and one attaining minimum reconstruction for every vector. Each reconstructed retained matrix equals \(W\) after permitted source permutation and rescaling. Restoring fixed one-hot measurement-error columns and inversely rescaling or permuting source laws preserves \(P_{\mathrm{obs}}\). Thus unequal image vectors have faithful, minimum-edge, observationally equivalent attaining countermodels.

The qualifier is relative Zariski generic in each fixed real canonical coefficient stratum. Its deleted proper algebraic sets are Lebesgue-null in the stratum's free Euclidean coordinates. Separability concerns source laws and remains separate. An exactly encoded rational input is deterministic, so no probability claim is made for rational points.

%%% FIELD proof-joint-column-singleton-criterion
The fixed-mleaf branch is \(\{e_u\}\). Suppose \(g_{\mathcal G,\ell}(u)=g\in[k]\). Every \(c\in L_u\) occurs as a legal pivot with center \(u\), so \(W_{uc}\ne0\). If \(L_u=\varnothing\), nonemptiness of \(P_g\) forces an attaining pair with center different from \(u\), whence \(z_u=1\) and \(S_{\bullet u}=\{e_u\}\).

If \(z_u=1\), then \(e_u\in S_{\bullet u}\). A normalized column equals \(e_u\) exactly when, for every \(v\ne u\),
\[
\frac{W_{vc}}{W_{uc}}=0\quad\Longleftrightarrow\quad W_{vc}=0.
\]
Thus the joint set is a singleton exactly under the stated zero tests, and its value is \(e_u\).

If \(z_u=0\) and \(L_u\ne\varnothing\), no \(e_u\) is adjoined, so singletonness is equality of all normalized column vectors. For \(c,d\in L_u\),
\[
\frac{W_{vc}}{W_{uc}}=\frac{W_{vd}}{W_{ud}}
\quad\Longleftrightarrow\quad
W_{vc}W_{ud}-W_{vd}W_{uc}=0,
\]
because both denominators are nonzero. Requiring this for every \(v\in V\) is equivalent to vector equality. These cases are exhaustive and prove both directions and the common values.

%%% FIELD proof-effect-column-bound-sharp
Fix \(q\ge1\). We construct a two-row, one-group model and give an explicit source-law proof of separability. Write
\[
W=(w_1,\ldots,w_q),\qquad w_j=\begin{pmatrix}a_j\\b_j\end{pmatrix},
\]
where \(a_jb_j\ne0\) and
\[
a_jb_h-a_hb_j\ne0\qquad(j\ne h).
\]
Thus the retained columns are pairwise nonparallel and none is a coordinate vector.

These inequalities can be imposed with Yang rank faithfulness at a rational point. The rational reference columns \((1,j)^{\mathsf T}\) satisfy the coordinate and minor inequalities and give the legal one-group canonical stratum verified below. On that nonempty stratum, \(\Theta_{2,q,\ell}^{\mathrm{gen}}\) deletes finitely many proper algebraic loci for rank recovery and simultaneous DOG faithfulness. After clearing pivot denominators, multiply their nonzero defining polynomials by the coordinate and minor polynomials to obtain a nonzero polynomial \(F\). Rational points are Zariski dense: a nonzero univariate polynomial has finitely many roots, and induction on the number of variables first chooses rational values where a nonzero coefficient polynomial survives and then a final rational coordinate outside the finite root set. Hence a rational \(W\) with \(F\ne0\) exists. By \(\text{lem:yang-algorithm-equivalence}\), its recovered group is \(R_1=\{u,v\}\), \(K_1=C_{\mathrm{src}}=[q]\); by \(\text{lem:yang-dog-faithfulness-preservation}\), Assumption 3 holds simultaneously throughout its finite DOG class.

To prove nonempty separability rather than assume it, let the \(q\) structural sources and two one-hot measurement-error sources be mutually independent Bernoulli variables, each charging both endpoints. The convex hull of the observed support is, up to translation,
\[
Z=\sum_{j=1}^{q}[0,w_j]+[0,e_u]+[0,e_v].
\]
Indeed every endpoint combination has positive probability, and convex hull commutes with a finite Minkowski sum.

The unoriented edge directions of this planar zonotope are exactly its generator lines. For a generator \(z\), choose a normal \(n\ne0\) with \(n^{\mathsf T}z=0\) and \(n^{\mathsf T}z'\ne0\) for every nonparallel generator \(z'\). The face maximizing \(n^{\mathsf T}x\) fixes an endpoint of every \([0,z']\) but leaves all of \([0,z]\) free, producing an edge parallel to \(z\). Conversely an exposed edge of a Minkowski sum of segments is parallel to a generator orthogonal to its exposing normal. Since the lines of \(w_1,\ldots,w_q,e_u,e_v\) are distinct, \(P_{\mathrm{obs}}\) determines them from its support hull. Known measurement labels identify the two coordinate lines; removing them leaves exactly the columns of \(W\) up to permutation and nonzero scaling. Thus this explicit independent non-Gaussian source law satisfies \(\text{ass:separability}\).

Select center \(r\in\{u,v\}\) and source \(c\in[q]\), and let \(l\) be the other row. Then
\[
Q_\omega=(1),\qquad A_\omega=(0),\qquad D_\omega=\frac{W_{lc}}{W_{rc}}\ne0.
\]
For each \(h\ne c\),
\[
(B^C_\omega)_h=W_{rh}\ne0,
\qquad
(B^L_\omega)_h
=\frac{W_{lh}W_{rc}-W_{lc}W_{rh}}{W_{rc}}\ne0.
\]
Thus every one of the \(2q\) choices is canonical. Each complementary latent source has the two children \(r,l\); for \(q=1\) there is no complementary latent. The exact edge count is
\[
E(\omega)=0+(q-1)+1+(q-1)=2q-1.
\]
There are no other assignments, so all are global minimizers.

For minimality, each root latent has mleaf child \(l\) and other child \(r\). The edge \(r\to l\) gives
\[
r\in\operatorname{An}(l),\qquad r\notin\operatorname{An}(r),
\]
so \(\operatorname{An}(l)\nsubseteq\operatorname{An}(r)\), excluding the witness in \(\text{lem:yang-minimality-criterion}\). For \(q=1\) the criterion is vacuous. One-hot measurement errors supply the known measured labels without changing the retained matrix.

Center \(u\) yields the \(q\) pairwise distinct vectors
\[
\frac{W_{\bullet j}}{W_{uj}}=\begin{pmatrix}1\\b_j/a_j\end{pmatrix},
\]
and center \(v\) makes \(u\) an mleaf, yielding \(e_u=(1,0)^{\mathsf T}\). Since every \(b_j/a_j\ne0\), the latter differs from all normalized columns. Therefore
\[
|S_{\bullet u}(W,\mathcal G(W,\ell),\ell)|=q+1,
\]
proving sharpness on the published class for every \(q\ge1\).

%%% FIELD statement-effect-column-bound-sharp
For every integer \(q\ge1\), the bound \(q+1\) in \(\text{thm:deterministic-joint-effect-image}\) is attained by a separable canonical, minimal, LV-SEM-ME rank-faithful Yang model with known labels: there are two measured substantive labels \(u\ne v\), one recovered group \(R_1=\{u,v\}\), \(K_1=C_{\mathrm{src}}=[q]\), and an exactly encoded rational retained mixing representative for which
\[
|S_{\bullet u}(W,\mathcal G(W,\ell),\ell)|=q+1.
\]
Every one of the \(2q\) legal row--column choices is a global minimum and has \(2q-1\) structural edges.

%%% FIELD proof-comp-edge-promise-np
An assignment \(\omega=((r_g,c_g))_{g=1}^{k}\) has at most \(k\le\min(p,q)\) pairs, so its binary encoding is polynomial in the input bit length. Given \(\omega\), check \(r_g\in R_g\), \(c_g\in K_g\), the optional equality \(\omega_g=(r,c)\), distinctness of selected rows and columns, known-label restrictions, and all explicitly encoded Algorithm-2 pivot and canonicality constraints. Reject any zero selected pivot.

Clear rational denominators and compute \(Q_\omega^{-1}\), \(A_\omega\), \(B^C_\omega\), \(D_\omega\), and \(B^L_\omega\) by fraction-free elimination and exact multiplication. Bit sizes remain polynomial. If an \(m\times m\) integer matrix has entries of at most \(h\) bits, the Leibniz bound gives determinant bit length
\[
O\bigl(m(h+\log m)\bigr).
\]
Clearing denominators raises \(h\) by at most the total input bit length. Cramer's rule bounds inverse numerators and denominators by such minors, and polynomially many subsequent rational operations preserve polynomial bit length.

Verify lower triangularity or acyclicity, latent-child, mleaf, label, and canonical support conditions by exact zero tests. The four blocks contain
\[
k^2+k(q-k)+(p-k)k+(p-k)(q-k)=pq
\]
positions, so support counting and comparison with \(t\) are polynomial. This verifier accepts exactly the promised yes instances, proving both versions are in promise-\(\mathrm{NP}\).

The optimum is an integer in \([0,pq]\), so binary search with an \(\mathrm{NP}\) oracle determines it in \(O(\log(pq+1))\) queries. For search, ask the \(\mathrm{NP}\) language whether an accepted optimal certificate extends each successive candidate bit prefix. One branch remains feasible until a full certificate is recovered. Thus optimization and witness recovery belong to \(\mathrm{FP}^{\mathrm{NP}}\). This gives neither a deterministic polynomial optimizer nor an unrestricted hardness result.

%%% FIELD proof-block-completion-decomposition
Fix a legal choice \(\omega\). Since every selected pair lies in one block, reordering selected rows and columns by block makes \(\Delta_\omega\) block diagonal. The cross-block zero hypothesis also makes \(\widetilde W[R(\omega),C(\omega)]\) block diagonal. Therefore
\[
Q_\omega
=\widetilde W[R(\omega),C(\omega)]\Delta_\omega^{-1}
=\bigoplus_{j=1}^{s}Q_{\omega_j},
\qquad
Q_\omega^{-1}=\bigoplus_{j=1}^{s}Q_{\omega_j}^{-1}.
\]
The matrices \(\widetilde W[R,H]\), \(\widetilde W[L,C]\), and \(\widetilde W[L,H]\) are block diagonal under the same order. Substitution into \(\text{def:normalized-reconstruction}\) gives
\[
A_\omega=\bigoplus_jA_{\omega_j},\qquad
B^C_\omega=\bigoplus_jB^C_{\omega_j},\qquad
D_\omega=\bigoplus_jD_{\omega_j},\qquad
B^L_\omega=\bigoplus_jB^L_{\omega_j}.
\]

Candidate membership, nonzero pivots, observation labels, mleaf restrictions, and latent-child conditions are local to a group and therefore to its block. A block direct sum is acyclic exactly when each block is acyclic, and no cross-block coefficient can violate canonicality. Hence a global assignment is legal exactly when all block restrictions are legal:
\[
\Omega=\prod_{j=1}^{s}\Omega_j.
\]
The support size of a direct sum is additive, so
\[
E(\omega)=\sum_{j=1}^{s}E_j(\omega_j).
\]
Minimizing over the Cartesian product gives
\[
e^\star
=\min_{(\omega_1,\ldots,\omega_s)\in\prod_j\Omega_j}\sum_jE_j(\omega_j)
=\sum_{j=1}^{s}\min_{\omega_j\in\Omega_j}E_j(\omega_j)
=\sum_{j=1}^{s}e_j^\star.
\]
Combining local minimizers attains equality. Conversely, replacing any nonminimal component of a purported global minimizer would strictly reduce the sum. Thus
\[
\Omega^\star=\prod_{j=1}^{s}\Omega_j^\star.
\]
A prescribed pair lies in a unique block and restricts only that factor.

%%% FIELD statement-block-interaction-tractability
Fix a nonnegative integer \(b\). Consider a promised rational instance admitting the block partitions in \(\text{lem:block-completion-decomposition}\). Call a group \(g\) nontrivial when \(|R_g||K_g|>1\), and suppose every block contains at most \(b\) nontrivial groups. Then both \(\operatorname{CompOpt}(\widetilde W,\mathcal G,\ell)\) and every prescribed-pair problem \(\operatorname{CompOpt}(\widetilde W,\mathcal G,\ell;g,r,c)\) are solvable deterministically in
\[
(pq)^b\operatorname{poly}(p,q,\text{input bit length})
\]
time. The algorithm returns an attaining legal completion whenever one exists and an exact infeasibility certificate otherwise. Arbitrarily many nontrivial groups are allowed overall. Every complete block assignment, including the unique choices of groups that are not nontrivial, is subjected to the exact \(\mathrm{COMP\text{-}EDGE}\) verification procedure; singleton candidate sets are not presumed legal.

%%% FIELD proof-block-interaction-tractability
Fix \(b\) independently of the input. In block \(j\), let \(I_j\) be its set of nontrivial groups and put \(m_g=|R_g||K_g|\). By hypothesis \(|I_j|\le b\), and \(m_g\le pq\). A group outside \(I_j\) contributes exactly one syntactic candidate pair. Hence the number of complete syntactic assignments in block \(j\) is
\[
\prod_{g\in I_j}m_g\le(pq)^{|I_j|}\le(pq)^b.
\]
This is only a candidate count, not a legality conclusion.

Enumerate these complete block assignments. For each, run the deterministic exact verification procedure constructed in \(\text{thm:comp-edge-promise-np}\), restricted to the block matrices. In particular, it checks the unique pair of every group with \(m_g=1\), rejects a zero pivot or any failed label, acyclicity, latent-child, mleaf, or canonical support condition, computes the four normalized blocks by fraction-free elimination, and counts exact support. The verifier proof bounds each run by a polynomial in \(p\), \(q\), and the input bit length. Retain a legal block assignment of minimum support, or declare the block infeasible if none is accepted.

By \(\text{lem:block-completion-decomposition}\), global legality is blockwise legality and
\[
E(\omega)=\sum_jE_j(\omega_j),\qquad e^\star=\sum_je_j^\star.
\]
If every block is feasible, combining retained local minimizers is a legal global minimizer. If one block is infeasible, the product formula for \(\Omega\) certifies global infeasibility. There are at most \(k\le\min(p,q)\) nonempty blocks containing candidate groups, so their number is absorbed into the polynomial factor. The total running time is
\[
(pq)^b\operatorname{poly}(p,q,\text{input bit length}).
\]

For a prescribed pair \((g,r,c)\), first reject if \((r,c)\notin R_g\times K_g\). Otherwise restrict enumeration only in the unique block containing \(g\) to assignments with \(\omega_g=(r,c)\); all other blocks are enumerated and verified as above. The count cannot increase, and the decomposition proves that the combined result is optimal under the prescription or that it is infeasible. Because the hypothesis is per block, arbitrarily many nontrivial groups may occur over arbitrarily many blocks. The faithful plateau is consistent with this theorem: all three of its nontrivial groups lie in one interacting block, so it is handled by joint block enumeration, not by a false additive per-group score.

%%% FIELD proof-plateau-certificate
Order the substantive variables as \((c_0,c_1,c_2,l_0,l_1,l_2)\) and retained sources as \((N_0,N_1,N_2,H_0,H_1)\). Consider
\[
C=\Gamma C+B^CH+N,\qquad L=DC+B^LH,
\]
where
\[
\Gamma=\begin{pmatrix}0&0&0\\14&0&0\\14&11&0\end{pmatrix},
\quad B^C=\begin{pmatrix}4&0\\10&16\\0&13\end{pmatrix},
\]
\[
D=\begin{pmatrix}4&0&0\\9&18&0\\0&0&7\end{pmatrix},
\quad B^L=\begin{pmatrix}0&0\\0&0\\5&0\end{pmatrix}.
\]
Since \(\Gamma\) is strictly lower triangular,
\[
Q=(I-\Gamma)^{-1}=\begin{pmatrix}1&0&0\\14&1&0\\168&11&1\end{pmatrix}.
\]
Exact multiplication gives
\[
QB^C=\begin{pmatrix}4&0\\66&16\\782&189\end{pmatrix},
\quad DQ=\begin{pmatrix}4&0&0\\261&18&0\\1176&77&7\end{pmatrix},
\]
\[
DQB^C+B^L=\begin{pmatrix}16&0\\1224&288\\5479&1323\end{pmatrix}.
\]
Thus the retained mixing matrix is the matrix in \(\mathcal I_{\mathrm{plat}}\).

Its twelve edges are
\[
c_0\to c_1, c_0\to c_2, c_1\to c_2, c_0\to l_0, c_0\to l_1, c_1\to l_1, c_2\to l_2,
\]
\[
H_0\to c_0, H_0\to c_1, H_0\to l_2, H_1\to c_1, H_1\to c_2.
\]
They are acyclic in the displayed causal order. The latent roots have respectively three and two children; each \(l_i\) has no substantive child and no distinct structural noise. Independent one-hot measurement errors on the six measured rows establish the remaining canonical observation equations without changing \(W\).

Rank faithfulness is checked against the corrected Assumption 3. Reorder retained columns as \((H_0,H_1,N_0,N_1,N_2)\):
\[
\widehat W=
\begin{pmatrix}
4&0&1&0&0\\66&16&14&1&0\\782&189&168&11&1\\
16&0&4&0&0\\1224&288&261&18&0\\5479&1323&1176&77&7
\end{pmatrix}.
\]
The strict ancestor and possible-parent sets are
\[
\begin{array}{c|c|c}
x&\operatorname{An}(x)&\operatorname{PP}(x)\\ \hline
c_0&\{H_0\}&\varnothing\\
c_1&\{H_0,H_1,c_0\}&\{c_0,l_0\}\\
c_2&\{H_0,H_1,c_0,c_1\}&\{c_0,c_1,l_0,l_1\}\\
l_0&\{H_0,c_0\}&\{c_0\}\\
l_1&\{H_0,H_1,c_0,c_1\}&\{c_0,c_1,l_0\}\\
l_2&\{H_0,H_1,c_0,c_1,c_2\}&\{c_0,c_1,c_2,l_0,l_1\}.
\end{array}
\]
For every Assumption 3 pair \((J,K)\), fraction-free integer elimination computes \(\operatorname{rank}(\widehat W[K\cup\{x\},J])\). Independently, enumeration of the \(2^8\) subsets of the eight structural vertices, allowing endpoints in the deleted set, gives the minimum vertex bottleneck from \(J\) to \(K\cup\{x\}\). Grouping the exhaustive equal values \(s\) gives
\[
\begin{array}{c|r|rrrrr}
x&\#\text{ tests}&s=0&s=1&s=2&s=3&s=4\\ \hline
c_0&2&1&1&0&0&0\\
c_1&32&4&16&12&0&0\\
c_2&256&16&78&122&40&0\\
l_0&8&2&6&0&0&0\\
l_1&128&8&50&70&0&0\\
l_2&1024&32&221&428&303&40.
\end{array}
\]
No pair mismatches, and
\[
2+32+256+8+128+1024=1450.
\]
For an mleaf-parent clause, \(\operatorname{PP}(V_j)\subseteq\operatorname{PP}(x)\) and \(\operatorname{An}(x)\setminus\{V_j\}\subseteq\operatorname{An}(x)\), so those pairs are included. This proves the exact rank-equals-bottleneck condition for both specified families.

The hypotheses of \(\text{lem:yang-minimality-criterion}\) now hold. The latent \(H_1\) has no mleaf child. The only mleaf child of \(H_0\) is \(l_2\), while \(c_0\) is another child and
\[
c_2\in\operatorname{An}(l_2),\qquad c_2\notin\operatorname{An}(c_0).
\]
Therefore \(\operatorname{An}(l_2)\nsubseteq\operatorname{An}(c_0)\), excluding the criterion's obstruction. This proves minimality in the corrected Assumption 4 sense.

For the completion enumeration, let \(a_i=0\) select \(c_i\) and \(a_i=1\) select \(l_i\), for \(i=0,1\); let \(z\) similarly select \(c_2\) or \(l_2\). Let \(x=1\) replace \(N_0\) by \(H_0\), and \(y=1\) replace \(N_1\) by \(H_1\). The third noise is \(N_2\). All possible pivots are nonzero:
\[
\begin{array}{c|cc}&N_0&H_0\\ \hline c_0&1&4\\l_0&4&16\end{array},
\quad
\begin{array}{c|cc}&N_1&H_1\\ \hline c_1&1&16\\l_1&18&288\end{array},
\quad W_{c_2N_2}=1,quad W_{l_2N_2}=7.
\]
The zeros above the group diagonal make every \(Q_\omega\) unit lower triangular. Exact substitution in the four reconstruction formulas gives, independently of \((a_0,a_1)\),
\[
\begin{array}{c|c|c}
(x,y,z)&(\|A\|_0,\|B^C\|_0,\|D\|_0,\|B^L\|_0)&\text{latent-column child counts}\\ \hline
(0,0,0)&(3,4,4,1)&(3,2)\\
(0,0,1)&(3,5,4,1)&(4,2)\\
(0,1,0)&(3,5,4,1)&(2,4)\\
(0,1,1)&(3,5,4,1)&(2,4)\\
(1,0,0)&(3,4,5,1)&(3,2)\\
(1,0,1)&(3,5,5,1)&(4,2)\\
(1,1,0)&(3,5,5,1)&(4,2)\\
(1,1,1)&(3,5,5,1)&(4,2).
\end{array}
\]
Thus all \(2^5=32\) assignments are legal, and summing supports yields
\[
E(a_0,a_1,z,x,y)=12+x+z+(1-z)y.
\]
Hence \(e^\star=12\), with equality exactly when \(x=y=z=0\). The four values of \((a_0,a_1)\) are precisely the stated minimizers.

At each of these four rational minimizers, the same exact rank--bottleneck enumeration after the corresponding two center transpositions again has \(1450\) tests, aggregate counts \((63,372,632,343,40)\) for ranks \((0,1,2,3,4)\), and no mismatch. Their graphs are the displayed graph with the selected transpositions \(c_i\leftrightarrow l_i\), \(i=0,1\). In each graph \(H_1\) has no mleaf child, while the only mleaf child \(l_2\) of \(H_0\) has ancestor \(c_2\), which is not an ancestor of the other selected child. Thus the same exact minimality criterion applies to every minimizer. The generic preservation lemma is therefore not being used illicitly to infer pointwise faithfulness of these rational countermodels.

Finally, \(((c_0,N_0),(c_1,H_1),(l_2,N_2))\) has \((x,y,z)=(0,1,1)\) and thirteen edges. Changing only group zero leaves at least thirteen; changing only group one still has \(z=1\) and leaves at least thirteen; changing only group two sets \(z=0\) while \(y=1\) and again leaves thirteen. Changing the second noise and third center together sets \(y=z=0\), and with \(x=0\) attains twelve. This proves the plateau and the coupled repair.
